\section{Rechnungsgrundlagen}\label{sec:rechnungsgrundlagen}

Rechnungsgrundlagen im Sinne des \S\,2 Abs.\,1 KVAV sind der
Rechnungszins, die Ausscheideordnung, die Kopfschäden, der
Sicherheitszuschlag und die sonstigen Zuschläge. Alle rechnungsmäßigen
Ansätze werden gemäß \S\,9 KVAV in überprüfbarer Weise belegt; die
Rechnungsgrundlagen sind gemäß \S\,2 Abs.\,3 KVAV mit ausreichenden
Sicherheiten versehen. Das Schlussalter aller Tafeln beträgt
$\omega = 110$.

\subsection{Rechnungszinssatz}\label{sec:zins}

Der Rechnungszins $i$ wird je Beobachtungseinheit festgelegt und
überschreitet den Höchstrechnungszins von 3{,}5\,\% gemäß \S\,4 VAG
nicht. Die Angemessenheit wird jährlich anhand des aktuariellen
Unternehmenszinses (AUZ) gemäß der DAV-Richtlinie
\glqq Der aktuarielle Unternehmenszins in der privaten
Krankenversicherung (AUZ)\grqq{} überprüft; die Herleitung des AUZ ist
in der Rahmen-TB \mbox{RTB-AUZ} (Version 2025-1) dokumentiert
(Abschnitt~\ref{sec:rahmentb}). Der AUZ des Unternehmens zum
Kalkulationsstichtag wird zusammen mit den rechnungsmäßigen Zinssätzen
in Anlage~\ref{sec:anlage21} ausgewiesen; die Zinssätze
liegen jeweils unterhalb des AUZ und sind damit vorsichtig gewählt. Es gilt der Diskontierungsfaktor
\begin{equation}
v = \frac{1}{1+i},
\label{eq:diskont}
\end{equation}
wobei $i$ den Rechnungszins der jeweiligen Beobachtungseinheit
bezeichnet. Eine unterjährige Verzinsung wird nicht angesetzt.

\subsection{Ausscheideordnung}\label{sec:ausscheideordnung}

Die Ausscheideordnung wird gemäß \S\,5 KVAV aus Sterbe- und
Stornowahrscheinlichkeiten zusammengesetzt; beide werden gesondert
nachgewiesen.

\subsubsection{Sterbewahrscheinlichkeiten}\label{sec:sterblichkeit}

Als rechnungsmäßige Sterbewahrscheinlichkeiten $q_x$ werden die Werte
der jeweils aktuellsten von der BaFin veröffentlichten
\textbf{PKV-Sterbetafel} verwendet, derzeit der PKV-Sterbetafel~2025
(Spalten $q_x$ für Männer und $q_y$ für Frauen, Alter 0 bis 102), ohne
Zu- oder Abschläge. Eine Prüfung ergab keine Anhaltspunkte dafür, dass
die Tafel für den Bestand des Tarifs ungeeignet ist.

\textbf{Fortschreibung oberhalb des Tafelendalters (Festlegung):} Die
PKV-Sterbetafel endet beim Alter 102. Für die Alter 103 bis 109 werden
die Sterbewahrscheinlichkeiten linear interpoliert zwischen $q_{102}$
und $q_{110} := 1$:
\begin{equation}
q_x = q_{102} + \bigl(1 - q_{102}\bigr)\cdot\frac{x-102}{110-102},
\qquad x = 103,\dots,109,
\label{eq:qfort}
\end{equation}
wobei $q_{102}$ den letzten Tafelwert bezeichnet. Im Schlussalter
$\omega = 110$ scheiden alle Versicherten rechnungsmäßig aus.

\subsubsection{Stornowahrscheinlichkeiten}\label{sec:storno}

Die rechnungsmäßigen Stornowahrscheinlichkeiten $w_x$ beruhen auf der
unternehmenseigenen Stornobeobachtung der Krankheitskostenvollversicherung
der Geschäftsjahre 2021 bis 2023 (Personenzählung unter Elimination von
zwei Beginnjahren zur Ausschaltung des Frühstornos; Methodik gemäß der
DAV-Richtlinie \glqq Festlegung von Stornotafeln in der privaten
Krankenversicherung\grqq). Die Werte werden geschlechtsunabhängig
angesetzt, da die Beobachtung keine signifikanten Unterschiede zwischen
Männern und Frauen ergab. Gegenüber den bisher angesetzten Werten
(Änderungsnr.~16) wurden die Stornowahrscheinlichkeiten in den Altern
21 bis 35 um durchschnittlich 0{,}5 Prozentpunkte gesenkt. Im
Schlussalter gilt $w_{110} := 0$, zusammen mit $q_{110} = 1$ also
$(q+w)_{110} = 1$.

\subsubsection{Zusammengesetzte Ausscheideordnung}\label{sec:zao}

Die Ausscheidewahrscheinlichkeit ist
\begin{equation}
(q+w)_x = q_x + w_x,
\label{eq:qw}
\end{equation}
wobei $q_x$ die Sterbe- und $w_x$ die Stornowahrscheinlichkeit des
Alters $x$ bezeichnet. Die \textbf{diskontierte
Übergangswahrscheinlichkeit} ist
\begin{equation}
DQ_x = v\cdot\bigl(1-(q+w)_x\bigr),
\label{eq:dq}
\end{equation}
wobei $v$ der Diskontierungsfaktor gemäß Gleichung~\eqref{eq:diskont}
ist. Zur Veranschaulichung ergibt sich hieraus die Bestandsfunktion
$l_{x+1} = \bigl(1-(q+w)_x\bigr)\cdot l_x$ mit dem Startwert
$l_{21} = 1\,000\,000$ sowie die diskontierten Lebenden
$D_x = v^x\cdot l_x$ mit $DQ_x = D_{x+1}/D_x$; für die Berechnungen
dieser TB werden ausschließlich die Rekursionen des
Abschnitts~\ref{sec:grundwerte} verwendet. Die vollständigen Tafeln
$q_x$, $w_x$, $(q+w)_x$ und $DQ_x$ enthält Anlage~\ref{sec:anlage221}.

\subsection{Kopfschaden}\label{sec:kopfschaden}

Die folgenden Unterabschnitte dokumentieren die Grundsätze der
Kopfschadenermittlung, die Umlage der Schwangerschafts- und
Mutterschaftskosten, die Herleitung der rechnungsmäßigen Profile und
Grundkopfschäden sowie die Kopfschäden der Kinder und Jugendlichen.

\subsubsection{Grundsätze}\label{sec:kopfschaden-grundsaetze}

Die Kalkulation erfolgt mit Profilen und Grundkopfschäden nach der
Methode von Rusam: Der rechnungsmäßige Kopfschaden des Alters $x$ ist
\begin{equation}
K_x = k_x \cdot G,
\label{eq:rusam}
\end{equation}
wobei $k_x$ den normierten Kopfschaden (Profilwert) des Alters $x$ und
$G$ den Grundkopfschaden der Beobachtungseinheit bezeichnet. Das
Normierungsalter ist 40, d.\,h. $k_{40} = 1$ und $G = K_{40}$.

Da der eigene Bestand des Tarifs für eine tarifindividuelle Herleitung
nicht ausreichend groß ist (\S\,6 Abs.\,3 KVAV), werden als Stützdaten
die von der BaFin veröffentlichten Kopfschadentafeln
(\textbf{PKV-Kopfschadenstatistik, Beobachtungsjahr 2022}) verwendet.
Das Leistungsversprechen des Tarifs entspricht in Leistungsumfang und
Selbstbehaltsstufe den folgenden Tafeln für Nichtbeihilfeberechtigte
(Versicherungsart~N):

\begin{center}\small
\begin{tabular}{@{}lll@{}}
\toprule
Leistungsbereich & BaFin-Tafel (Männer/Frauen/S-Kosten) & Bemerkung \\
\midrule
ambulant & 01\_KKV\_amb\_N\_\{M,W,S\}\_0-100 & Selbstbehalt 0--100 EUR \\
stationär & 02\_KKV\_stat\_N\_\{M,W,S\}\_MB & Mehrbettzimmer \\
Zahnbehandlung & 03\_KKV\_zahn\_N\_\{M,W,S\}\_Beh & \\
Zahnersatz & 03\_KKV\_zahn\_N\_\{M,W,S\}\_Ers\_70-80 & Erstattung 70--80\,\% \\
\bottomrule
\end{tabular}
\end{center}

\subsubsection{Umlage der Schwangerschafts- und Mutterschaftskosten}
\label{sec:sumlage}

Die Tafeln der Frauen enthalten die Schwangerschafts- und
Mutterschaftskosten (S-Kosten). Diese werden seit 2009 (AGG)
geschlechtsunabhängig umgelegt. Je Leistungsbereich und Alter $x$ wird
die Umlage
\begin{equation}
u_x = \frac{n^{W}_x \cdot K^{S}_x}{n^{M}_x + n^{W}_x}
\label{eq:sumlage}
\end{equation}
gebildet, wobei $K^{S}_x$ die beobachteten S-Kosten je versicherte Frau
(Tafeln mit Kennung~S), $n^{M}_x$ und $n^{W}_x$ die Bestände
(Personenjahre) der Männer- bzw. Frauentafel des Leistungsbereichs
bezeichnen. Die beobachteten Teilkopfschäden nach Umlage sind
\begin{equation}
K^{\mathrm{LB},M}_x = K^{\mathrm{LB}}_x(\text{M}) + u_x,
\qquad
K^{\mathrm{LB},W}_x = K^{\mathrm{LB}}_x(\text{W}) - K^{S}_x + u_x,
\label{eq:teilkopf}
\end{equation}
wobei $K^{\mathrm{LB}}_x(\cdot)$ den Kopfschaden der jeweiligen
BaFin-Tafel bezeichnet. In Altern ohne S-Kosten (keine
Tafelwerte) ist $K^S_x = 0$.

\subsubsection{Herleitung der rechnungsmäßigen Profile und
Grundkopfschäden}\label{sec:kopfschaden-herleitung}

Der beobachtete Gesamtkopfschaden je Geschlecht ist die Summe der
Teilkopfschäden über die vier Leistungsbereiche:
\begin{equation}
K^{\mathrm{beob}}_x
= K^{\mathrm{amb}}_x + K^{\mathrm{stat}}_x + K^{\mathrm{zahnB}}_x
+ K^{\mathrm{zahnE}}_x ,
\qquad x = 0,\dots,100.
\label{eq:ksumme}
\end{equation}
Die vollständige Herleitung je Alter enthält Anlage~\ref{sec:anlage232}.

\textbf{Fortschreibung oberhalb des Tafelendalters (Festlegung):} Die
BaFin-Tafeln enden beim Alter 100. Für die Alter 101 bis 110 wird der
Gesamtkopfschaden geometrisch fortgeschrieben:
\begin{equation}
K^{\mathrm{beob}}_x = K^{\mathrm{beob}}_{100}\cdot g^{\,x-100},
\qquad
g = \Bigl(K^{\mathrm{beob}}_{100} / K^{\mathrm{beob}}_{95}\Bigr)^{1/5},
\qquad x = 101,\dots,110,
\label{eq:kfort}
\end{equation}
wobei $g$ die mittlere jährliche Veränderungsrate der Alter 95 bis 100
bezeichnet; die festgesetzten Veränderungsraten je Beobachtungseinheit
werden in Anlage~\ref{sec:anlage232} ausgewiesen.

Die Profilwerte ergeben sich durch Normierung auf das Alter 40,
\begin{equation}
k_x = \frac{K^{\mathrm{beob}}_x}{K^{\mathrm{beob}}_{40}},
\label{eq:profil}
\end{equation}
und werden auf vier Nachkommastellen gerundet.
\textbf{Monotonisierung (Festlegung):} Ab dem Monotonisierungsalter
$\tilde{x}$ werden die Profilwerte monoton nicht fallend angesetzt,
\begin{equation}
k_x := \max\{k_{x-1},\, k_x\}, \qquad x = \tilde{x}+1,\dots,110,
\label{eq:mono}
\end{equation}
da die beobachteten Profile in den schwach besetzten hohen Altern
fallen; der festgesetzte Wert $\tilde{x}$ wird in
Anlage~\ref{sec:anlage233} ausgewiesen. Die Monotonisierung wirkt als
Sicherheitsmarge und stellt
zusammen mit der Beitragsformel des Abschnitts~\ref{sec:beitragsraten}
monoton nicht fallende Beiträge sicher (\S\,146 Abs.\,2 VAG).

Der \textbf{rechnungsmäßige Grundkopfschaden} wird je
Beobachtungseinheit als $G = K^{\mathrm{beob}}_{40}$ festgesetzt:
Die festgesetzten Werte $G^{M}$ und $G^{F}$, die rechnungsmäßigen
Profile $k_x$ und die rechnungsmäßigen Kopfschäden $K_x = k_x\cdot G$
werden -- da sie sich mit jeder Beitragsanpassung ändern können --
in Anlage~\ref{sec:anlage233} ausgewiesen; dort ist auch
der Nachweis der ausreichenden Sicherheit der Grundkopfschäden
dokumentiert.

\subsubsection{Kopfschäden der Kinder und Jugendlichen}
\label{sec:kopfschaden-kinder}

Für die Beobachtungseinheit der Kinder (0--15) wird ein
altersunabhängiger, geschlechtsunabhängiger mittlerer Kopfschaden
angesetzt; für die Jugendlichen (16--20) werden altersunabhängige,
nach Geschlecht getrennte mittlere Kopfschäden der männlichen und der
weiblichen Jugendlichen angesetzt (jeweils bestandsgewichtetes Mittel
der Gesamtkopfschäden nach S-Umlage über die Altersgruppe, Gewichte:
Bestände der Ambulanttafeln). Die festgesetzten Werte
$\bar{K}_{\mathrm{Ki}}$, $\bar{K}^{M}_{\mathrm{Ju}}$ und
$\bar{K}^{F}_{\mathrm{Ju}}$ werden in Anlage~\ref{sec:anlage233}
ausgewiesen.

\subsection{Zuschläge}\label{sec:zuschlaege}

Die Zuschläge umfassen den Sicherheitszuschlag, die Kostenzuschläge,
den Zuschlag für die erfolgsunabhängige Beitragsrückerstattung, den
Zuschlag für Standard- und Basistarif sowie die Zillmerung.

\subsubsection{Sicherheitszuschlag}\label{sec:sicherheitszuschlag}

Der Sicherheitszuschlag $\sigma$ gemäß \S\,7 KVAV wird einheitlich
für alle Beobachtungseinheiten und Alter als Prozentsatz des
Bruttobeitrags angesetzt und beträgt mindestens den gesetzlichen
Mindestsatz von 5~Prozent; der festgesetzte Wert wird in
Anlage~\ref{sec:anlage241} ausgewiesen. Er ist
nicht bereits in anderen Rechnungsgrundlagen enthalten.
\textbf{Begründung der Höhe:} Der gesetzliche Mindestsatz wird als
ausreichend angesehen, da die übrigen Rechnungsgrundlagen bereits
eigene Sicherheiten enthalten (vorsichtiger Rechnungszins unterhalb
des AUZ, Sicherheitsmarge der Grundkopfschäden gemäß
Abschnitt~\ref{sec:kopfschaden-herleitung}, Monotonisierung der
Profile, Sterbetafel ohne Abschläge). Rundungen der
Beitragsberechnung gehen zu Gunsten oder zu Lasten des
Sicherheitszuschlags (Abschnitt~\ref{sec:rundung}).

\subsubsection{Kostenzuschläge}\label{sec:kosten}

Die Kostensätze werden in der Kostenrechnung des Unternehmens je
Beobachtungseinheit hergeleitet (Rahmen-TB \mbox{RTB-Kosten}) und als
Prozentsätze des Bruttobeitrags angesetzt. Die einzelnen Kostenarten
(\S\,8 Abs.\,1 KVAV) -- unmittelbare Abschlusskosten (laufender Zuschlag
$\alpha_u$), mittelbare Abschlusskosten $\alpha_m$,
Schadenregulierungskosten $\rho$ und sonstige Verwaltungskosten $\beta$
-- und ihre Summen $\Delta_x$ werden zusammen mit der einbezogenen
Version der Rahmen-TB in Anlage~\ref{sec:anlage241} ausgewiesen, da sie
sich mit jeder Beitragsanpassung ändern können.

Der laufende Zuschlag für unmittelbare Abschlusskosten entfällt nach
Vollendung des 65.~Lebensjahres (\S\,8 Abs.\,3 Satz~3 KVAV).
Gemäß \S\,8 Abs.\,4 KVAV dürfen in die Prämien -- mit Ausnahme der
Zillmerung und der Zuschläge nach \S\,8 Abs.\,1 Nr.\,6 und 8 KVAV --
nur altersunabhängige absolute Kostenzuschläge eingerechnet werden.
Die prozentualen Kostensätze werden daher in altersunabhängige
Zuschläge $\gamma$ (EUR p.\,a.) umgerechnet. Die Umrechnung erfolgt
über das \textbf{Gamma-Alter} $x_\gamma$, das auf ganze Jahre
gerundete bestandsgewichtete Durchschnittsalter der
Beobachtungseinheit in der Krankheitskostenvollversicherung des
Unternehmens zum 31.12.2023 (Bestandsstatistik; fiktive Festlegung:
Männer $x_\gamma = \GammaAlterM$, Frauen $x_\gamma = \GammaAlterF$):
\begin{equation}
\gamma^{\mathrm{Kosten}}_{\le 64}
= \frac{12\,\Delta_{\le 64}\cdot\bigl(G\cdot A_{x_\gamma}
  + \gamma^{\mathrm{euBR}}\cdot a_{x_\gamma}\bigr)}
       {12\,\bigl(1-\Delta_{\le 64}-\sigma-\Omega_{\le 64}\bigr)\cdot
        a_{x_\gamma} - \alpha_{x_\gamma}},
\qquad
\gamma^{\mathrm{Kosten}}_{\ge 65}
= \gamma^{\mathrm{Kosten}}_{\le 64}\cdot
  \frac{\Delta_{\ge 65}}{\Delta_{\le 64}},
\label{eq:gammakosten}
\end{equation}
wobei $A_{x_\gamma}$ und $a_{x_\gamma}$ die Barwerte des Gamma-Alters
gemäß Abschnitt~\ref{sec:grundwerte}, $\sigma$ den Sicherheitszuschlag
gemäß Abschnitt~\ref{sec:sicherheitszuschlag},
$\Omega$ den Zuschlag gemäß Abschnitt~\ref{sec:omega} und
$\alpha_{x_\gamma}$ den Zillmersatz gemäß
Abschnitt~\ref{sec:zillmerung} bezeichnen. Der Ausdruck
$12\,\Delta\cdot b_{x_\gamma}$ entspricht dem jährlichen prozentualen
Kostenbetrag der Musterperson des Gamma-Alters.

\subsubsection{Zuschlag für die erfolgsunabhängige
Beitragsrückerstattung}\label{sec:eubr}

Der Zuschlag für die erfolgsunabhängige Beitragsrückerstattung
(\S\,8 Abs.\,1 Nr.\,5 KVAV) wird als absoluter Zuschlag angesetzt
(Festlegung auf Basis der Auszahlungsquoten der Geschäftsjahre
2021--2023). Die festgesetzten Werte $\gamma^{\mathrm{euBR}}$ je
Beobachtungseinheit werden in Anlage~\ref{sec:anlage241}
ausgewiesen.

\subsubsection{Zuschlag für Standardtarif und Basistarif}
\label{sec:omega}

Der beitragsproportionale Zuschlag zur Umlage der Begrenzung der
Beitragshöhe im Basistarif (\S\,154 VAG, \S\,8 Abs.\,1 Nr.\,6 KVAV)
und für den Standardtarif (\S\,8 Abs.\,1 Nr.\,8 KVAV) beträgt
\begin{equation}
\Omega_x = \begin{cases}
\OmegaU\,\% & x \le 64,\\[2pt]
\OmegaO\,\% & x \ge 65,
\end{cases}
\label{eq:omegaz}
\end{equation}
da die Standardtarif-Komponente nach Vollendung des 65.~Lebensjahres
entfällt (\S\,8 Abs.\,5 KVAV). \textbf{Zusammensetzung:} Die
Basistarif-Komponente (0{,}6\,\%, branchenweit einheitlich vom
PKV-Verband ermittelt) und die Standardtarif-Komponente (0{,}1\,\%,
$x \le 64$) sind poolbedingt; ihre Herleitung ist in den Technischen
Berechnungsgrundlagen des Basistarifs bzw. Standardtarifs des
Unternehmens dokumentiert. Ein unternehmensindividueller Anteil wird
nicht erhoben (Festlegung).

\subsubsection{Summe der Zuschläge}\label{sec:zuschlagsumme}

Die altersunabhängigen Zuschläge je Beobachtungseinheit ergeben sich
als Summe
\begin{equation}
\gamma_x = \gamma^{\mathrm{euBR}} + \gamma^{\mathrm{Kosten}}_x;
\label{eq:gammasumme}
\end{equation}
die festgesetzten Werte (einschließlich der Kinder und Jugendlichen)
sowie die Summe der beitragsproportionalen Zuschläge
$\sigma + \Omega_x$ werden in Anlage~\ref{sec:anlage241}
ausgewiesen. Eine Übersicht aller Zuschläge enthält Anlage~\ref{sec:anlage241}.

\subsubsection{Zillmerung}\label{sec:zillmerung}

Zur nachgelagerten Finanzierung der unmittelbaren Abschlusskosten wird
gezillmert. Der Zillmersatz $\alpha_x$ in Vielfachen des
Monatsbeitrags beträgt
\begin{equation}
\alpha_x = \begin{cases}
0 & x \le 20,\\
\AlphaZ & 21 \le x \le 75,\\
3{,}5 - 0{,}7\cdot(x-75) & 76 \le x \le 79,\\
0 & x \ge 80.
\end{cases}
\label{eq:zillmer}
\end{equation}
Die Zillmerkosten eines Vertrags betragen $\alpha_{x_e}\cdot b_{x_e}$
(einmalig bei Vertragsbeginn), wobei $x_e$ das Eintrittsalter und
$b_{x_e}$ den Monatsbeitrag bezeichnet. Der Nachweis der zulässigen
Höhe des Zillmersatzes gemäß \S\,8 Abs.\,3 KVAV wird in
Abschnitt~\ref{sec:bilanzar} und Anlage~\ref{sec:anlage522} geführt.
